% !TEX program = xelatex
% Lernplatt - by Can Siebert
% Kurs 22 (Bo 143) - Tag 7 - Loesungsheft
% Plattform-SEO, Content, Retention & Bewertungsmanagement

\documentclass[11pt,a4paper,oneside]{article}

\usepackage{polyglossia}
\setdefaultlanguage{german}

\usepackage[a4paper,margin=2.5cm]{geometry}

\usepackage{fontspec}
\defaultfontfeatures{Ligatures=TeX}

\IfFontExistsTF{Charter}{\setmainfont{Charter}}{%
  \setmainfont{Noto Serif}}
\IfFontExistsTF{Source Sans 3}{\setsansfont{Source Sans 3}}{%
  \IfFontExistsTF{Source Sans Pro}{\setsansfont{Source Sans Pro}}{%
    \setsansfont{Liberation Sans}}}

\newcommand{\caveatfontfallback}{\itshape}
\IfFontExistsTF{Caveat}{\newfontfamily\caveatfont{Caveat}[Scale=1.15]}{\newcommand\caveatfont{\caveatfontfallback}}

\setmonofont{Liberation Mono}

\usepackage{microtype}
\usepackage{eurosym}
\linespread{1.15}

\usepackage{xcolor}
\definecolor{lpInk}{RGB}{20,28,38}
\definecolor{lpAccent}{RGB}{22,90,140}
\definecolor{lpAccentLight}{RGB}{210,228,240}
\definecolor{lpAmber}{RGB}{222,138,30}
\definecolor{lpAmberLight}{RGB}{255,243,217}
\definecolor{lpAmberBorder}{RGB}{196,118,18}
\definecolor{lpGray}{RGB}{96,104,114}
\definecolor{lpGrayLight}{RGB}{238,240,243}
\definecolor{lpGreen}{RGB}{34,120,72}
\definecolor{lpGreenLight}{RGB}{222,238,228}
\definecolor{lpGreenBorder}{RGB}{24,96,58}
\definecolor{lpL1}{RGB}{34,120,72}
\definecolor{lpL2}{RGB}{22,90,140}
\definecolor{lpL3}{RGB}{170,42,52}

\usepackage[most]{tcolorbox}
\tcbuselibrary{breakable,skins}

\newtcolorbox{infobox}[1][Hinweis]{enhanced, breakable, colback=lpAccentLight, colframe=lpAccent, boxrule=0.6pt, arc=2pt, left=10pt,right=10pt,top=8pt,bottom=8pt, fonttitle=\sffamily\bfseries\color{white}, coltitle=white, title={#1}, attach boxed title to top left={xshift=8pt,yshift=-8pt}, boxed title style={size=small, colback=lpAccent, colframe=lpAccent, boxrule=0.4pt, arc=1pt}, top=14pt}

\newtcolorbox{loesungbox}[1][Musterlösung]{enhanced, breakable, colback=lpGreenLight, colframe=lpGreenBorder, boxrule=0.6pt, arc=2pt, left=10pt,right=10pt,top=8pt,bottom=8pt, fonttitle=\sffamily\bfseries\color{white}, coltitle=white, title={#1}, attach boxed title to top left={xshift=8pt,yshift=-8pt}, boxed title style={size=small, colback=lpGreenBorder, colframe=lpGreenBorder, boxrule=0.4pt, arc=1pt}, top=14pt}

\newtcolorbox{merkbox}[1][Managementhinweis]{enhanced, breakable, colback=lpAmberLight, colframe=lpAmberBorder, boxrule=0.6pt, arc=2pt, left=10pt,right=10pt,top=8pt,bottom=8pt, fonttitle=\sffamily\bfseries\color{white}, coltitle=white, title={#1}, attach boxed title to top left={xshift=8pt,yshift=-8pt}, boxed title style={size=small, colback=lpAmberBorder, colframe=lpAmberBorder, boxrule=0.4pt, arc=1pt}, top=14pt}

\usepackage{enumitem}
\setlist{itemsep=2pt, topsep=4pt, parsep=0pt}
\setlist[itemize,1]{label={\textbullet},leftmargin=1.6em}
\setlist[itemize,2]{label={\textendash},leftmargin=1.4em}
\setlist[enumerate,1]{label=\arabic*., leftmargin=1.8em}

\usepackage{tabularx}
\usepackage{array}
\usepackage{booktabs}
\usepackage{longtable}

\usepackage{fancyhdr}
\usepackage{lastpage}
\pagestyle{fancy}
\fancyhf{}
\renewcommand{\headrulewidth}{0.3pt}
\renewcommand{\footrulewidth}{0pt}
\fancyhead[L]{\footnotesize\sffamily\color{lpGray}Lösungsheft \textperiodcentered{} Kurs 22 \textperiodcentered{} Tag 7}
\fancyhead[R]{\footnotesize\sffamily\color{lpGray}Plattform-SEO, Retention \& Bewertungen}
\fancyfoot[L]{\footnotesize\sffamily\color{lpGray}\textcopyright{} Can Siebert}
\fancyfoot[R]{\footnotesize\sffamily\color{lpGray}\thepage{} / \pageref{LastPage}}

\fancypagestyle{plain}{\fancyhf{}\renewcommand{\headrulewidth}{0pt}\fancyfoot[C]{\footnotesize\sffamily\color{lpGray}\thepage{} / \pageref{LastPage}}}

\usepackage[explicit]{titlesec}
\titleformat{\section}[block]{\sffamily\Large\bfseries\color{lpAccent}}{\thesection.}{0.6em}{#1}[{\vspace{2pt}\color{lpAccent}\titlerule[0.6pt]}]
\titleformat{\subsection}[block]{\sffamily\large\bfseries\color{lpInk}}{\thesubsection}{0.6em}{#1}
\titlespacing*{\section}{0pt}{14pt}{8pt}
\titlespacing*{\subsection}{0pt}{10pt}{4pt}

\usepackage[hidelinks,unicode]{hyperref}

\newcommand{\akzent}[1]{{\caveatfont\color{lpAmber}#1}}
\newcommand{\fachbegriff}[1]{\textbf{#1}}

\newcommand{\niveaubadge}[2]{\tcbox[on line, colback=#1, colframe=#1, coltext=white, boxrule=0pt, arc=2pt, left=5pt,right=5pt,top=1pt,bottom=1pt, nobeforeafter]{\sffamily\bfseries\footnotesize #2}}
\newcommand{\nivLi}{\niveaubadge{lpL1}{L1\,\textbullet\,Wissen}}
\newcommand{\nivLii}{\niveaubadge{lpL2}{L2\,\textbullet\,Anwendung}}
\newcommand{\nivLiii}{\niveaubadge{lpL3}{L3\,\textbullet\,Management}}

\newcommand{\zeit}[1]{\tcbox[on line, colback=lpGrayLight, colframe=lpGrayLight, coltext=lpInk, boxrule=0pt, arc=2pt, left=5pt,right=5pt,top=1pt,bottom=1pt, nobeforeafter]{\sffamily\footnotesize\textbf{$\sim$\,#1}}}

\newcommand{\aufgabenkopf}[4]{\par\vspace{6pt}\noindent{\sffamily\Large\bfseries\color{lpAccent}Aufgabe #1 \textendash{} #2}\par\vspace{2pt}\noindent{\color{lpAccent}\rule{\linewidth}{0.6pt}}\par\vspace{4pt}\noindent #3 \hfill \zeit{#4}\par\vspace{6pt}}

\newcommand{\ytquery}[1]{\par\vspace{4pt}\noindent{\footnotesize\sffamily\color{lpGray}\textbf{YouTube-Suchquery:} \texttt{#1}}\par}

\begin{document}

\begin{titlepage}
  \pagestyle{empty}
  \vspace*{1.5cm}
  \noindent{\sffamily\footnotesize\color{lpGray}LERNPLATT \textperiodcentered{} BY CAN SIEBERT}\\[2pt]
  \noindent{\color{lpAccent}\rule{\linewidth}{1.2pt}}\\[4pt]
  \noindent{\sffamily\footnotesize\color{lpGray}LÖSUNGSHEFT \textperiodcentered{} MODUL 4 \textperiodcentered{} TAG 7 \textperiodcentered{} KURS 22 (Bo 143) \textperiodcentered{} 8.\,Juni 2026}

  \vspace{3cm}
  {\sffamily\fontsize{30}{34}\selectfont\bfseries\color{lpInk}Lösungsheft\\Tag 7\par}

  \vspace{0.7cm}
  {\sffamily\large\color{lpGray}Musterlösungen zu den elf Aufgaben\\des Aufgabenhefts \textendash{} Plattform-SEO,\\Retention \& Bewertungsmanagement.\par}

  \vspace{1.5cm}
  {\caveatfont\LARGE\color{lpGreen}Musterlösungen \textperiodcentered{} nicht die einzige Antwort}\par

  \vfill

  \noindent\begin{minipage}{0.55\linewidth}
    {\sffamily\footnotesize\color{lpGray}Hinweis:}\\
    {\sffamily\small\color{lpInk}Musterlösungen sind Diskussionsbasis.}\\[10pt]
    {\sffamily\footnotesize\color{lpGray}Begleitend:}\\
    {\sffamily\small\color{lpInk}Aufgabenheft \& Lehrskript Tag 7}
  \end{minipage}\hfill
  \begin{minipage}{0.4\linewidth}
    \raggedleft
    {\sffamily\footnotesize\color{lpGray}Dozent:}\\
    {\sffamily\small\color{lpInk}Can Siebert}\\[10pt]
    {\sffamily\footnotesize\color{lpGray}Niveau-Mix:}\\
    {\sffamily\small\color{lpInk}3\,L1 \textperiodcentered{} 5\,L2 \textperiodcentered{} 3\,L3}
  \end{minipage}

  \vspace{0.5cm}
  \noindent{\color{lpAccent}\rule{\linewidth}{0.6pt}}
\end{titlepage}

\section*{Hinweise zu den Musterlösungen}
\thispagestyle{plain}

\noindent Eine mögliche Lösung. Mehrere Begründungswege sind legitim, solange sie:

\begin{itemize}
  \item den Begriffsapparat von Tag 7 nutzen (Ranking-Faktoren, Keyword-Arten, Bewertungsmanagement, Retention),
  \item Annahmen sichtbar machen,
  \item Zielkonflikte (Sichtbarkeit vs.\ Marke, Klicks vs.\ Marge) benennen,
  \item zu einer klaren Entscheidung führen.
\end{itemize}

\clearpage

\section*{Niveau L1 \textendash{} Wissen \& Reproduktion}
\addcontentsline{toc}{section}{L1 -- Wissen}

% --- AUFGABE 1 -----------------------------------------------------------
% LERNPLATT_HILFSVIDEOS
\section*{Hilfsvideos zu diesem Tag}
\addcontentsline{toc}{section}{Hilfsvideos zu diesem Tag}
\thispagestyle{plain}

\noindent Die folgenden YouTube-Erkl\"arvideos vermitteln die Inhalte, die Sie zur Bearbeitung der Aufgaben ben\"otigen. Klicken Sie auf den Titel, um das Video direkt zu \"offnen; die vollst\"andige URL steht in Klammern.

\begin{description}[style=nextline,leftmargin=18pt,labelindent=0pt,itemsep=8pt]
  \item[1.~Amazon-Listing optimieren — Plattform-SEO in der Praxis] \href{https://youtu.be/sU6n178V0-I}{\textcolor{lpAccent}{https://youtu.be/sU6n178V0-I}}\\[2pt]
  {\footnotesize\color{lpGray}(URL: \nolinkurl{https://youtu.be/sU6n178V0-I})}
  \item[2.~Longtail-Keywords — die Geheimwaffe für qualifizierte Anfragen] \href{https://youtu.be/8i51lMBCnko}{\textcolor{lpAccent}{https://youtu.be/8i51lMBCnko}}\\[2pt]
  {\footnotesize\color{lpGray}(URL: \nolinkurl{https://youtu.be/8i51lMBCnko})}
  \item[3.~SEO-Grundlagen — Tutorial für die strategische Sicht] \href{https://youtu.be/NbG9qXoO5N0}{\textcolor{lpAccent}{https://youtu.be/NbG9qXoO5N0}}\\[2pt]
  {\footnotesize\color{lpGray}(URL: \nolinkurl{https://youtu.be/NbG9qXoO5N0})}
  \item[4.~Kundenbindung und Wiederkauf — vom Erstkauf zum CLV] \href{https://youtu.be/WJHr8R1s4y8}{\textcolor{lpAccent}{https://youtu.be/WJHr8R1s4y8}}\\[2pt]
  {\footnotesize\color{lpGray}(URL: \nolinkurl{https://youtu.be/WJHr8R1s4y8})}
\end{description}
\vspace{0.6em}
\noindent{\footnotesize\color{lpGray}\textit{Tipp:} \"Offnen Sie das Video, lassen Sie es im Hintergrund laufen und arbeiten Sie parallel an der Aufgabe.}

\clearpage

\aufgabenkopf{1}{Ranking-Faktoren auf Plattformen}{\nivLi}{6\,min}

\ytquery{Plattform SEO Ranking Faktoren Conversion Bewertung Produkttitel}

\begin{loesungbox}
\textbf{1. Zuordnung:}

\smallskip
\begin{tabularx}{\linewidth}{@{}lX@{}}
\toprule
\textbf{Kategorie} & \textbf{Faktor} \\
\midrule
Content-Faktoren                & Suchbegriff-Relevanz, Produkttitel, Beschreibung, Kategorie, Bilder \\
\midrule
Performance-Faktoren            & Lieferfähigkeit, Klickrate, Conversion Rate, Retourenquote, Bewertungen, Verkäuferperformance \\
\bottomrule
\end{tabularx}

\textbf{2. Wirkt auf Sichtbarkeit \emph{und} Vertrauen:}
\emph{Bewertungen.} Sie sind im Algorithmus ein Ranking-Signal und beim Kunden ein Vertrauenssignal. Bewertungen verbessern beides gleichzeitig \textendash{} aber schlechte Bewertungen schädigen auch beides gleichzeitig.

\textbf{3. Häufig vernachlässigt im Mittelstand:}
\emph{Bilder} (Anbieter nutzen meist nur ein Produktfoto, keine Anwendungssituationen) und \emph{Retourenquote} (wird oft nicht aktiv gesteuert, ist aber für viele Plattform-Algorithmen ein zentrales Verkäufer-Performance-Signal).
\end{loesungbox}

\clearpage

% --- AUFGABE 2 -----------------------------------------------------------
\aufgabenkopf{2}{Fünf Keyword-Arten unterscheiden}{\nivLi}{6\,min}

\ytquery{Keyword Arten Hauptkeyword Long Tail kaufnah problemorientiert SEO}

\begin{loesungbox}
\textbf{1. Definitionen:}
\begin{itemize}
  \item \emph{Hauptkeyword:} kurzes, generisches Suchwort mit hohem Volumen. Beispiel: ``Akkuschrauber''.
  \item \emph{Nebenkeyword:} sekundärer, oft verwandter Begriff. Beispiel: ``Bohrschrauber''.
  \item \emph{Long-Tail-Keyword:} mehrteilig, spezifisch. Beispiel: ``Akkuschrauber mit zwei Akkus für Handwerk''.
  \item \emph{Problemorientiert:} beschreibt ein Kundenproblem. Beispiel: ``Akkuschrauber für Dauereinsatz lange Akkulaufzeit''.
  \item \emph{Kaufnah:} hohe Kaufabsicht im Suchbegriff. Beispiel: ``Akkuschrauber Profi 18V Set kaufen''.
\end{itemize}

\textbf{2. Volumen vs.\ Conversion:}
\begin{itemize}
  \item Höchstes Volumen, niedrigste Conversion: \emph{Hauptkeyword} (zu unspezifisch).
  \item Niedrigstes Volumen, höchste Conversion: \emph{Kaufnah} oder spezielles \emph{Long-Tail} (hohe Kaufabsicht, geringe Streuverluste).
\end{itemize}

\textbf{3. Wichtig für erklärungsbedürftige B2B-Produkte:}
\emph{Problemorientierte Keywords} (treffen den echten Bedarf) und \emph{Long-Tail-Keywords} (treffen die spezifische Anwendung). Hauptkeywords liefern dort Streuverluste, weil B2B-Kunden präziser suchen.
\end{loesungbox}

\clearpage

% --- AUFGABE 3 -----------------------------------------------------------
\aufgabenkopf{3}{Sichtbarkeit \textendash{} Klicks \textendash{} Conversion \textendash{} Profitabilität}{\nivLi}{8\,min}

\ytquery{Plattform SEO Sichtbarkeit Klicks Conversion Profitabilität KPI}

\begin{loesungbox}
\textbf{1. Vier Ebenen:}
\begin{itemize}
  \item \emph{Sichtbarkeit:} Wie oft taucht das Produkt in Suchergebnissen / Kategorien auf (Impressionen, Ranking-Position).
  \item \emph{Klicks:} Wie viele Nutzer klicken auf das Produkt (Klickrate).
  \item \emph{Conversion:} Wie viele Klicker:innen kaufen (Conversion Rate).
  \item \emph{Profitabilität:} Wie viel Deckungsbeitrag pro Verkauf nach Plattformgebühren und Logistik.
\end{itemize}

\textbf{2. Warum mehr Sichtbarkeit weniger Profitabilität bedeuten kann:}
Wenn Sichtbarkeit über \emph{falsche} Keywords erzeugt wird (z.\,B.\ ``billig''), kommen falsche Kunden mit falschen Erwartungen. Klickrate steigt, aber Conversion sinkt, Retourenquote steigt, Bewertungen verschlechtern sich. Das Plattform-Ranking sinkt langfristig \textendash{} obwohl es kurzfristig durch mehr Aktivität gestiegen ist.

\textbf{3. Trennkennzahlen:}
\begin{itemize}
  \item Klicks vs.\ Conversion: \emph{Conversion Rate} (wie viele Klicker:innen kaufen).
  \item Conversion vs.\ Profitabilität: \emph{Deckungsbeitrag pro Verkauf} (Marge nach allen Plattformkosten).
\end{itemize}

\begin{merkbox}[Managementhinweis]
\emph{Plattform-SEO darf nie auf nur einer Ebene optimiert werden.} Wer nur Sichtbarkeit optimiert, kauft Klicks. Wer nur Klicks optimiert, kauft Conversion. Wer nur Conversion optimiert, kauft Marge. Die Architektur muss alle vier Ebenen ins Verhältnis setzen.
\end{merkbox}
\end{loesungbox}

\clearpage

\section*{Niveau L2 \textendash{} Anwendung \& Transfer}
\addcontentsline{toc}{section}{L2 -- Anwendung}

% --- AUFGABE 4 -----------------------------------------------------------
\aufgabenkopf{4}{CleanPro Supplies \textendash{} Plattformangebot analysieren}{\nivLii}{25\,min}

\ytquery{B2B Marktplatz SEO Industriesauger Produkttitel Beschreibung Optimierung}

\begin{loesungbox}
\textbf{1. Fünf Schwächen:}
\begin{itemize}
  \item Produkttitel ``Sauger Modell X200'' enthält keine relevanten Suchbegriffe (``industriell'', ``nass'', ``trocken'', ``gewerblich'').
  \item Beschreibung ist generisch (``gute Qualität''), kein Nutzenversprechen, keine Zielgruppe.
  \item Nur ein Produktfoto \textendash{} keine Anwendungssituationen.
  \item Keine Kompatibilitäts- / Zubehörhinweise.
  \item Wenige Bewertungen (18) im Vergleich zum Wettbewerb \textendash{} schwaches Vertrauenssignal.
  \item Keine FAQ.
  \item Keine sichtbaren Differenzierungsmerkmale (Lebensdauer, Saugleistung, Ersatzteilverfügbarkeit).
\end{itemize}

\textbf{2. Drei Suchbegriffe:}
\begin{itemize}
  \item ``Industriesauger nass trocken gewerblich''.
  \item ``Nass-Trockensauger Werkstatt langlebig''.
  \item ``Gewerbesauger Ersatzteile verfügbar''.
\end{itemize}

\textbf{3. Verbesserter Produkttitel:}
``Industrieller Nass-Trockensauger für Gewerbe \& Werkstatt \textendash{} robust, langlebig, mit verfügbaren Ersatzteilen (Modell X200, 1.600\,W)''. Begründung: enthält Hauptkeyword + Zielgruppe + Differenzierung + technische Information.

\textbf{4. Drei Inhalte in der Beschreibung:}
\begin{itemize}
  \item \emph{Anwendungsfälle:} ``für Gebäudereinigung, Werkstätten, Handwerk, Facility-Management'' \textendash{} signalisiert Zielgruppen-Fit.
  \item \emph{Technische Spezifikationen:} Saugleistung, Behältergröße, Stromaufnahme, Schlauchlänge, Filter.
  \item \emph{Differenzierung:} ``Lebensdauer 5+\,Jahre, Ersatzteile garantiert für 10\,Jahre verfügbar, robustes Edelstahlgehäuse'' \textendash{} adressiert den Premium-Charakter.
\end{itemize}

\textbf{5. Zwei Vertrauensmaßnahmen:}
\begin{itemize}
  \item Strukturiertes Bewertungsmanagement (aktiv um Bewertungen bitten, alle Bewertungen öffentlich beantworten).
  \item Anwendungsbilder (Sauger in echten Einsatzumgebungen \textendash{} Werkstatt, Baustelle, Gewerbe).
\end{itemize}

\begin{merkbox}[Managementhinweis]
\emph{Premiumprodukte gewinnen nicht über Preis, sondern über Sichtbarkeit der Differenzierung.} Wenn die Plattform-Liste fünf Sauger zeigt und nur einer die Lebensdauer und Ersatzteilverfügbarkeit nennt, gewinnt dieser \textendash{} selbst wenn er teurer ist.
\end{merkbox}
\end{loesungbox}

\clearpage

% --- AUFGABE 5 -----------------------------------------------------------
\aufgabenkopf{5}{CleanPro \textendash{} Keyword- und Bewertungsstrategie}{\nivLii}{30\,min}

\ytquery{Marktplatz SEO Keyword Strategie B2B Bewertungsmanagement Premium}

\begin{loesungbox}
\textbf{1. Keyword-Priorisierung:}
\begin{itemize}
  \item \emph{Priorität 1:} ``gewerblicher Nass Trockensauger'' (mittel/Volumen, gute Zielgruppe).
  \item \emph{Priorität 2:} ``Industriesauger nass trocken'' (hohes Volumen, hoher Wettbewerb, aber zur Marke passend).
  \item \emph{Priorität 3:} ``robuster Sauger Werkstatt'' (hoher Praxisbezug).
  \item \emph{Priorität 4:} ``Profi Reinigungssauger langlebig'' (Premium-Positionierung).
  \item \emph{Verworfen:} ``billiger Industriesauger'' \textendash{} markenwidrig, falsche Kundenerwartung.
\end{itemize}

\textbf{2. Verteilung Titel / Beschreibung / FAQ:}
\begin{itemize}
  \item \emph{Titel:} ``Industriesauger nass trocken'' + ``gewerblich'' (Hauptkeyword + Qualifier).
  \item \emph{Beschreibung:} alle vier priorisierten Keywords natürlich integriert; problemorientierte Phrasen (``für Dauereinsatz'', ``für Werkstatt'').
  \item \emph{FAQ:} ``Ist der Sauger für Wasser geeignet?'', ``Wie lange hält der Filter?'', ``Welche Ersatzteile sind verfügbar?''.
\end{itemize}

\textbf{3. Drei Content-Maßnahmen für Conversion:}
\begin{itemize}
  \item Vergleichstabelle mit zwei Alternativ-Modellen (zeigt die eigene Stärke ohne Aggression).
  \item Anwendungsbilder in mindestens drei Umgebungen (Werkstatt, Baustelle, Gewerbe).
  \item Detaillierte Ersatzteilliste mit Verfügbarkeitsgarantie.
\end{itemize}

\textbf{4. Drei Maßnahmen Bewertungsmanagement:}
\begin{itemize}
  \item Wöchentliche Auswertung neuer Bewertungen mit Themen-Clusterung (Lieferzeit, Produkt, Service).
  \item Antwort innerhalb 48\,h auf jede Bewertung (auch positive).
  \item Strukturierte Bitte um Bewertung 14\,--\,21\,Tage nach Erhalt (E-Mail, nicht Spam-aggressiv).
\end{itemize}

\textbf{5. ``Billiger Industriesauger'' \textendash{} ja oder nein?}
\textbf{Nein.} Dieses Keyword erzeugt zwar Sichtbarkeit, aber falsche Kundenerwartung. Kunden, die ``billig'' suchen und ein 600-\euro-Premium-Modell finden, fühlen sich getäuscht $\rightarrow$ schlechte Bewertungen, höhere Retourenquote, schlechteres Plattform-Ranking. Markenwirkung sinkt langfristig. Das vermeintliche Sichtbarkeitsgeschenk ist ein Reputationsproblem.

\begin{merkbox}[Lessons Learned]
\emph{Sichtbarkeit über falsche Keywords ist ein Boomerang.} Die schnellen Klicks entwickeln sich in den folgenden Monaten zu schlechten Bewertungen und zu sinkender Algorithmus-Performance. Faustregel: \emph{Wenn eine Keyword-Auswahl falsche Kundenerwartung erzeugt, ist sie strategisch falsch \textendash{} auch wenn sie kurzfristig Klicks bringt.}
\end{merkbox}
\end{loesungbox}

\clearpage

% --- AUFGABE 6 -----------------------------------------------------------
\aufgabenkopf{6}{Bewertungsmanagement als Frühwarnsystem}{\nivLii}{20\,min}

\ytquery{Bewertungsmanagement B2B Plattform 1 Stern Antwort Frühwarnsystem}

\begin{loesungbox}
\textbf{1. Bewertungen als Lernsystem:}
Bewertungen liefern strukturierte Rückmeldung von echten Kunden \emph{nach} der Kaufentscheidung. Sie zeigen Lücken in Produktbeschreibung, Logistik, Service, Erwartungsmanagement und Produktqualität. Wer Bewertungen systematisch auswertet, kann Probleme identifizieren, bevor sie sich im Umsatz niederschlagen \textendash{} Bewertungen sind also nicht nur ein Verkaufssignal, sondern auch ein Frühwarnsystem für die Organisation.

\textbf{2. Klassifizierung der vier Bewertungsthemen:}
\begin{itemize}
  \item ``Lieferung 14 statt 5 Tage'' \textendash{} \emph{operativ} (Logistik / Lieferanten).
  \item ``Beschreibung passte nicht'' \textendash{} \emph{operativ} (Content-Pflege), aber Hinweis auf strategisches Problem (\emph{wir wecken falsche Erwartungen}).
  \item ``Service hat sich nicht gemeldet'' \textendash{} \emph{operativ} (Service-SLA).
  \item ``Produkt fühlt sich nicht hochwertig an'' \textendash{} \emph{strategisch}, weil es die Markenpositionierung trifft.
\end{itemize}

\textbf{3. Drei Plattform-Dashboard-KPIs:}
\begin{itemize}
  \item Bewertungsdurchschnitt (Trend über Monate).
  \item Anteil 1-2-Stern-Bewertungen (Frühwarnsignal).
  \item Hauptkritikpunkte (Top-3 Themen-Cluster).
\end{itemize}

\textbf{4. Antwort auf 1-Stern-Bewertung (Vorlage):}
``Sehr geehrte:r [Name], vielen Dank, dass Sie uns Ihre Erfahrung mitteilen. Es tut uns leid, dass [konkrete Kritik aufgreifen]. Wir haben Ihren Fall an unseren Service weitergeleitet \textendash{} bitte melden Sie sich unter [Kontakt], damit wir [konkrete Lösung] anbieten können. Ihre Rückmeldung hilft uns, [Prozess / Produkt] zu verbessern. \\ Mit freundlichen Grüßen, [Name], CleanPro Supplies.''

\begin{merkbox}[Lessons Learned]
\emph{Eine professionell beantwortete 1-Stern-Bewertung wirkt oft stärker als eine unbeantwortete 5-Stern-Bewertung.} Andere Kunden lesen die Antworten und beurteilen daran die Service-Kultur des Anbieters.
\end{merkbox}
\end{loesungbox}

\clearpage

% --- AUFGABE 7 -----------------------------------------------------------
\aufgabenkopf{7}{ToolMaster Pro \textendash{} Plattform-Performance optimieren}{\nivLii}{30\,min}

\ytquery{ToolMaster B2B Plattform SEO Akkuschrauber Bewertungsmanagement Retention}

\begin{loesungbox}
\textbf{1. Ursachenanalyse:}
\begin{itemize}
  \item Produkte qualitativ stark, aber Nutzen auf der Plattform nicht sichtbar.
  \item Technische Produkttitel greifen nicht die Suchlogik der Kunden auf.
  \item Fehlende Kompatibilitätsinformationen erzeugen Unsicherheit + negative Bewertungen.
  \item Bilder + Beschreibungen unterstützen Kaufentscheidung nicht.
  \item Bewertungsmanagement reaktiv oder fehlt.
  \item Wiederkaufpotenziale bei Zubehör werden nicht genutzt.
\end{itemize}

\textbf{2. Keyword-Strategie für Akkuschrauber + Zubehör:}
\begin{itemize}
  \item \emph{Haupt:} ``Profi Akkuschrauber'', ``Akkuschrauber Handwerk'', ``Akkuschrauber Werkstatt''.
  \item \emph{Long-Tail:} ``Akkuschrauber mit Ersatzakku für Handwerker'', ``robuster Akkuschrauber für Montage'', ``Akkuschrauber Zubehör Set kompatibel''.
  \item \emph{Problemorientiert:} ``Akkuschrauber für Dauereinsatz'', ``Akkuschrauber mit langer Akkulaufzeit''.
  \item \emph{Bewusst nicht:} ``billig'', ``Schnäppchen'' \textendash{} markenwidrig.
\end{itemize}

\textbf{3. Optimierter Titel + Beschreibung-Struktur:}
\begin{itemize}
  \item \emph{Titel (vorher):} ``Akkuschrauber TX 18V Set''.
  \item \emph{Titel (nachher):} ``Profi Akkuschrauber 18V Set für Handwerk \& Montage \textendash{} mit 2 Akkus, Schnellladegerät und robustem Transportkoffer''.
  \item \emph{Beschreibung-Struktur:} Kurz-Nutzeneinstieg $\rightarrow$ Hauptvorteile (Leistung, Akkulaufzeit, Robustheit) $\rightarrow$ Technische Daten $\rightarrow$ Kompatibilität (Akkus, Bits, Zubehör) $\rightarrow$ Anwendungsfälle $\rightarrow$ Lieferumfang $\rightarrow$ FAQ.
\end{itemize}

\textbf{4. Drei Content-Maßnahmen für Vertrauen \& Conversion:}
\begin{itemize}
  \item Anwendungsbilder mit Kontext (Montage, Innenausbau, Werkstatt).
  \item Vergleichstabelle Basis-/Profi-/Zubehör-Set.
  \item Downloadbare Kompatibilitätsliste für Bits, Akkus, Zubehör.
\end{itemize}

\textbf{5. Drei Retention-Maßnahmen:}
\begin{itemize}
  \item Zubehörsets + Ersatzakkus als strukturierte Folgeangebote.
  \item Bundle-Angebote für professionelle Anwender (Akkuschrauber + 2 Akkus + Koffer + Zubehörset).
  \item Nachkaufimpulse über E-Mail (sofern Kundendaten verfügbar) oder Plattform-Empfehlungslogik.
\end{itemize}

\begin{merkbox}[Managementhinweis]
\emph{Zubehör- und Ersatzteile sind im B2B-Plattformgeschäft oft das Margen-Reservoir.} Wer nur das Hauptgerät optimiert, lässt 30\,--\,50\,\% des CLV liegen. Bundle-Logik schützt zusätzlich vor Preisvergleich.
\end{merkbox}
\end{loesungbox}

\clearpage

% --- AUFGABE 8 -----------------------------------------------------------
\aufgabenkopf{8}{Retention-Logik bei Verbrauchs- und Zubehörprodukten}{\nivLii}{15\,min}

\ytquery{B2B Plattform Retention Wiederkauf Verbrauchsmaterial Bundle CLV}

\begin{loesungbox}
\textbf{1. Warum Retention oft profitabler ist:}
Neukundengewinnung im Plattformkontext kostet pro qualifizierten Kunden 50\,--\,200\,\euro{} CAC (Sichtbarkeit, Werbung, Plattformgebühren). Bestandskunden haben einen niedrigeren CAC-Anteil pro Wiederkauf, höhere Conversion-Raten und weniger Retouren. Verbrauchs- und Zubehörprodukte werden zudem häufiger gekauft \textendash{} höhere Frequenz, planbare Umsätze.

\textbf{2. Drei Wiederkauf-Trigger für ToolMaster:}
\begin{itemize}
  \item \emph{Zeitbasiert:} 6\,Monate nach Erstkauf eines Akkuschraubers $\rightarrow$ Erinnerung an Ersatzakkus / Verschleißteile.
  \item \emph{Verbrauchsbasiert:} bei bekanntem Verbrauchsmuster (z.\,B.\ Bits-Set alle 4\,Monate) automatische Empfehlung.
  \item \emph{Anlassbasiert:} neue Produktversion / Zubehör-Erweiterung $\rightarrow$ gezielte E-Mail an Bestandskunden.
\end{itemize}

\textbf{3. Zwei Plattformfunktionen:}
\emph{Bundle-Logik} (z.\,B.\ Akkuschrauber + 2 Akkus + Werkzeugkoffer in einem Listing) und \emph{Verbrauchsmaterial-Subscription} (falls Plattform es unterstützt, z.\,B.\ ``Save \& Subscribe'' bei Amazon Business). Beide reduzieren Preisvergleich und schützen Marge.

\textbf{4. Retention-Kennzahl:}
\emph{Wiederkaufrate innerhalb 12\,Monaten} (Anteil der Kunden, die innerhalb eines Jahres nach Erstkauf erneut bei uns kaufen). Sekundär: durchschnittlicher \emph{Customer Lifetime Value} nach 24\,Monaten.
\end{loesungbox}

\clearpage

\section*{Niveau L3 \textendash{} Management \& Entscheidungsarchitektur}
\addcontentsline{toc}{section}{L3 -- Management}

% --- AUFGABE 9 -----------------------------------------------------------
\aufgabenkopf{9}{ToolMaster Pro \textendash{} 6-Monats-Optimierungsentscheidung}{\nivLiii}{45\,min}

\ytquery{Plattform Performance B2B Senior CRO Bewertungsmanagement Budget Marge}

\begin{loesungbox}
\textbf{1. Vier Hebel:}
\textbf{Keywords + Content} (Titel, Beschreibung, FAQ), \textbf{Bilder + Anwendungsvideos}, \textbf{Bewertungsmanagement + Retention-Logik}, \textbf{Bundle- und Wiederkauf-Mechaniken}.

\textbf{2. Budgetverteilung (90.000\,\euro):}
\begin{itemize}
  \item 25.000\,\euro{} \textendash{} Keyword-Analyse + Content-Optimierung (Titel, Beschreibung, FAQ).
  \item 20.000\,\euro{} \textendash{} Bilder + Anwendungsvideos + Vergleichstabellen.
  \item 15.000\,\euro{} \textendash{} FAQ + Kompatibilitätsdaten + Produktdaten-Qualität.
  \item 15.000\,\euro{} \textendash{} Bewertungsmanagement (Prozess, Reporting, Antwortressourcen).
  \item 10.000\,\euro{} \textendash{} Bundle- und Retention-Mechaniken.
  \item 5.000\,\euro{} \textendash{} Tests, KPI-Setup, A/B-Auswertungen.
\end{itemize}

\textbf{3. Bewertungsmanagement-System:}
\begin{itemize}
  \item \emph{Wer wertet aus:} Kundenservice + Produktmanagement gemeinsam, wöchentlich.
  \item \emph{Wer antwortet:} Kundenservice innerhalb 48\,h auf jede Bewertung.
  \item \emph{Eskalation:} bei wiederkehrender Kritik (3+ Bewertungen zum gleichen Thema) $\rightarrow$ Eskalation an Produktmanagement (Stufe 1), Vertriebsleitung (Stufe 2), GF (Stufe 3 bei kritischen Themen).
  \item \emph{Lernschleife:} Monatlich Bewertungs-Review-Meeting; Top-3 Themen werden zu Verbesserungsmaßnahmen.
\end{itemize}

\textbf{4. Entscheidung gegen Preissenkung:}
ToolMaster Pro verzichtet bewusst auf aggressive Preissenkungen, obwohl der Plattform-Algorithmus niedrige Preise belohnt. Aus Senior-Level-Perspektive: Die Marge (24\,\%) ist die Voraussetzung für Service, Qualität und Markenpflege. Eine 10\,\%-Preissenkung würde die Marge auf 14\,\% drücken \textendash{} und damit weder Service noch Qualitätserhalt finanzieren. Stattdessen wird in Sichtbarkeit der \emph{Qualitätsunterschiede} investiert (Content, Bilder, Bewertungen). Eine reaktive Preisstrategie wäre der erste Schritt in eine strukturelle Margenerosion, aus der sich ein Anbieter selten erholt.

\textbf{5. Vier KPIs:}
\begin{itemize}
  \item Conversion Rate (vs.\ Vorperiode).
  \item Bewertungsdurchschnitt + Anteil neuer 4-5-Stern-Bewertungen.
  \item Wiederkaufrate (12-Monats-Sicht) bei Bestandskunden.
  \item Deckungsbeitrag pro Plattform-Listing (nach Plattformgebühr).
\end{itemize}

\textbf{6. Entscheidung unter Unsicherheit:}
Investition von 20.000\,\euro{} in Anwendungsbilder und kurze Videos \emph{ohne} klaren ROI-Beweis. Annahme: \emph{Kunden, die Produkte in Anwendungssituationen sehen, kaufen mit höherer Wahrscheinlichkeit und mit weniger Retouren}. Lernindikator: nach 90\,Tagen testbarer Effekt auf Conversion Rate und Retourenquote bei 5 Pilot-Listings. Wenn keine Verbesserung $\geq$\,15\,\%: Re-Allokation in FAQ und Kompatibilitätsdaten.

\begin{merkbox}[Senior-Level-Hinweis]
\emph{Sichtbarkeit kann man kaufen \textendash{} Vertrauen muss man aufbauen.} Wer im Plattformkontext nur Sichtbarkeit kauft, optimiert Quartalszahlen. Wer Vertrauen aufbaut (Bewertungen, Service, Kompatibilität), schafft strukturellen Wettbewerbsvorteil.
\end{merkbox}
\end{loesungbox}

\clearpage

% --- AUFGABE 10 ----------------------------------------------------------
\aufgabenkopf{10}{Sichtbarkeit kaufen \textendash{} aber welche?}{\nivLiii}{30\,min}

\ytquery{Plattform Sichtbarkeit Sponsored Products Premium Listing Marke Strategie}

\begin{loesungbox}
\textbf{1. Klassifizierung:}

\smallskip
\begin{tabularx}{\linewidth}{@{}lX@{}}
\toprule
\textbf{Typ} & \textbf{Hebel} \\
\midrule
Aufbauend (langfristiger Wert)  & Organisches Ranking (Content, Bewertungen) \textperiodcentered{} Algorithmus-konformes Verhalten (Lieferzeit, Retourenquote) \\
\midrule
Kaufend (kurzfristiger Effekt)  & Premium-Listing \textperiodcentered{} Plattform-Werbung (Sponsored) \textperiodcentered{} Sale-Aktionen \\
\bottomrule
\end{tabularx}

\textbf{2. Für eine Qualitätsmarke strategisch wertvoller:}
\emph{Organisches Ranking} (Content + Bewertungen) und \emph{algorithmus-konformes Verhalten} (kurze Lieferzeit, niedrige Retourenquote). Beide bauen \emph{strukturelle} Sichtbarkeit auf, die nicht erodieren, wenn das Werbebudget gekürzt wird. Sie verstärken die Markenwahrnehmung; Premium-Listing und Sponsored Ads dagegen können sogar Marke beschädigen, wenn sie zu offensichtlich als ``gekaufte'' Sichtbarkeit wahrgenommen werden.

\textbf{3. Kurzfristig wirksam, strategisch riskant:}
\emph{Sale-Aktionen} (trainieren Kunden auf Rabatte und beschädigen Marken-Anker) und \emph{Premium-Listing} (kann als Manipulation des Bewertungssystems wahrgenommen werden \textendash{} und wird bei nachlassender Performance vom Algorithmus selbst zurückgenommen).

\textbf{4. Faustregel:}
\emph{Sichtbarkeit kaufen ist akzeptabel, um eine bereits aufgebaute Position zu verstärken \textendash{} aber nicht, um eine schwache Position zu kompensieren.} Wer auf gekaufter Sichtbarkeit aufbaut, ohne organische Substanz, baut auf Sand: Sobald die Werbung gestoppt wird, fällt das Listing unter den Tisch.

\textbf{5. Falsche Sichtbarkeit als späteres Reputationsproblem:}
Beispiel: Premium-Akkuschrauber für 449\,\euro{} wird über Keyword ``billiger Akkuschrauber'' beworben. Klickrate steigt kurzfristig (Sichtbarkeit). Kunden klicken, sehen den Preis, sind enttäuscht, viele kaufen nicht $\rightarrow$ Klickrate ohne Conversion lässt das Listing im Algorithmus abrutschen. Wenige, die kaufen, fühlen sich getäuscht $\rightarrow$ schlechte Bewertungen. Die Sichtbarkeit von gestern wird das Reputationsproblem von morgen.

\begin{merkbox}[Lessons Learned]
\emph{Plattform-Sichtbarkeit ist ein Geld-Marker, der die Substanz nicht ersetzt.} Wer organische Substanz hat (Content, Bewertungen, Reaktionszeit), darf Sichtbarkeit kaufen, um sie zu verstärken. Wer keine Substanz hat, kauft sich nur in die Sichtbarkeit eines schwachen Listings hinein \textendash{} und produziert Folgekosten.
\end{merkbox}
\end{loesungbox}

\clearpage

% --- AUFGABE 11 ----------------------------------------------------------
\aufgabenkopf{11}{Transferprojekt \textendash{} Plattform-Performance-Architektur}{\nivLiii}{45\,min}

\ytquery{Plattform Performance Architektur B2B CRO Bewertungsmanagement Marge}

\begin{loesungbox}
\emph{Musterlösung als Entscheidungsarchitektur \textendash{} andere Konfigurationen möglich.}

\smallskip
\textbf{1. Zielbild 12\,--\,18\,Monate:}
Das Unternehmen ist auf den priorisierten Plattformen als \emph{Qualitätsanbieter mit klarer Differenzierung} sichtbar, hat einen Bewertungsdurchschnitt $\geq$\,4,6 und mindestens dreimal so viele Bewertungen wie heute. Conversion Rate steigt um 30\,\%, Wiederkaufrate verdoppelt sich, Deckungsbeitrag pro Listing bleibt $\geq$\,18\,\%. Keine reaktiven Preisaktionen.

\textbf{2. Segmentierung Plattformkunden + Suchintentionen:}
\begin{itemize}
  \item Professionelle Anwender (Handwerk, Werkstatt) \textendash{} Suchintention: Qualität + Verlässlichkeit.
  \item Industrieller Bedarf \textendash{} Suchintention: Spezifikationen + Service.
  \item Wiederkaufkunden \textendash{} Suchintention: Ersatzteile + Kompatibilität.
\end{itemize}

\textbf{3. Keyword-Architektur:}
\begin{itemize}
  \item \emph{Hauptkeywords:} ``Profi [Produkt]'', ``[Produkt] für Handwerk''.
  \item \emph{Long-Tail:} produkt- und anwendungs-spezifische Kombinationen.
  \item \emph{Problemorientiert:} ``für Dauereinsatz'', ``mit langer Lebensdauer''.
  \item \emph{Bewusst ausgeschlossen:} ``billig'', ``Schnäppchen'', generische Wettbewerber-Vergleichs-Keywords.
\end{itemize}

\textbf{4. Content-Architektur:}
\begin{itemize}
  \item Titel: $\leq$\,150\,Zeichen, Hauptkeyword + Differenzierung.
  \item Beschreibung: Nutzeneinstieg $\rightarrow$ Vorteile $\rightarrow$ Spezifikationen $\rightarrow$ Kompatibilität $\rightarrow$ Anwendungen $\rightarrow$ Lieferumfang.
  \item Bilder: mind.\ 5 (Produkt + Anwendungen + Detail + Zubehör + Vergleich).
  \item Video: 30\,--\,60\,Sek.\ Anwendungsvideo (wenn Plattform unterstützt).
  \item FAQ: mind.\ 5 Antworten (Kompatibilität, Ersatzteile, Garantie, Service, Lieferung).
  \item Vergleichstabelle: zwei Alternativ-Modelle (z.\,B.\ Basis vs.\ Profi vs.\ Industrie).
\end{itemize}

\textbf{5. Retention-Logik:}
\begin{itemize}
  \item Wiederkauf-Trigger zeitbasiert + verbrauchsbasiert.
  \item Bundle-Angebote zur Vergleichbarkeitsreduktion.
  \item Verbrauchsmaterial-Subscription (sofern Plattform unterstützt).
  \item Bestandskunden-Pflege über E-Mail (wo Daten verfügbar) + Plattform-Empfehlungen.
\end{itemize}

\textbf{6. Bewertungsmanagement-System:}
\begin{itemize}
  \item Wöchentliche Auswertung; Themen-Clusterung.
  \item Antwort auf jede Bewertung in 48\,h.
  \item Eskalation: 3+ ähnliche Kritiken $\rightarrow$ Produktmanagement.
  \item Lernschleife: monatliches Review; Top-Themen werden Verbesserungs-Tickets.
  \item Aktive (regelkonforme) Bewertungs-Anfrage 14\,--\,21\,Tage nach Lieferung.
\end{itemize}

\textbf{7. KPI-Struktur:}
\begin{itemize}
  \item Operativ: Sichtbarkeit pro Listing, Klickrate, Conversion Rate, Retourenquote.
  \item Management: Bewertungsdurchschnitt, Anzahl neuer Bewertungen, Wiederkaufrate, Deckungsbeitrag.
  \item Strategisch: Marktanteil Qualitätssegment, CLV, Anteil organische vs.\ gekaufte Sichtbarkeit.
\end{itemize}

\textbf{8. Governance:}
\begin{itemize}
  \item CRO: Strategie + Budget.
  \item Head of Marketplace: operative Steuerung pro Plattform.
  \item Marketing: Content + Bilder.
  \item Produktmanagement: Spezifikationen + Vergleichsinhalte.
  \item Kundenservice: Bewertungsmanagement.
  \item Monatlicher Steuerungskreis Plattform-Performance.
\end{itemize}

\textbf{9. Risikoanalyse:}
\begin{itemize}
  \item Preisdruck $\rightarrow$ Bundle, Service-Mehrwert, Qualitätssichtbarkeit.
  \item Falsche Kundenerwartung $\rightarrow$ keine Preis-Keywords, klare Beschreibung.
  \item Schlechte Bewertungen $\rightarrow$ Bewertungsmanagement + Frühwarnsystem.
  \item Plattformabhängigkeit $\rightarrow$ Multi-Plattform-Präsenz + eigener Direktkanal.
  \item Margenverlust $\rightarrow$ Deckungsbeitrag-KPI als Steuerungs-KPI, nicht nur Umsatz.
\end{itemize}

\textbf{10. Drei Managemententscheidungen unter Unsicherheit:}
\begin{enumerate}
  \item Investition in Content/Bilder \emph{vor} sichtbarem Conversion-Effekt. Annahme: Kunden im Premium-Segment werten Content höher als Preis. Wenn falsch: Re-Allokation in Bewertungsmanagement.
  \item Verzicht auf Sponsored-Ads-Kampagne, obwohl der Plattform-Algorithmus sie belohnen würde. Annahme: organische Substanz schlägt gekaufte Sichtbarkeit langfristig.
  \item Bundle-Logik als Schutz vor Preisvergleich \emph{bevor} Bundle-Akzeptanz nachgewiesen ist. Annahme: B2B-Kunden honorieren Komplettlösungen.
\end{enumerate}

\textbf{Zusatz \textendash{} Verzicht auf kurzfristige Reichweite:}
Bewusster Verzicht auf eine 30-Tage-Rabattaktion über alle Plattformen. Aus Senior-Level-Sicht: Eine Rabattaktion zerstört den Preisanker für Bestandskunden, signalisiert ``hier wird man billig'' und führt zu Folge-Aktionen, weil das Listing ohne Rabatt schlechter performt. Stattdessen wird das Budget in Content und Bewertungen investiert.

\begin{merkbox}[Senior-Level-Hinweis]
\emph{Plattform-Performance ist eine Verteidigung gegen den Preisdruck-Algorithmus.} Wer dem Algorithmus folgt (``mehr Sichtbarkeit, niedrigere Preise''), wird zum Commodity-Anbieter. Wer ihm widersteht (Qualität sichtbar machen), bleibt eine Marke.
\end{merkbox}
\end{loesungbox}

\clearpage

\section*{Abschluss}
\thispagestyle{plain}

\noindent Diskussionsbasis. Vergleichen Sie Ihre Annahmen, Ihre Zielkonflikte und Ihre Entscheidungen.

\medskip
\begin{infobox}[Nächster Schritt]
Coaching-Frage: Welche Keyword-/Content-/Bewertungs-Entscheidung haben Sie aus Markenüberlegung getroffen \textendash{} und würden Sie diese auch dann verteidigen, wenn der Plattform-Algorithmus Ihr Listing kurzfristig schlechter rankt?
\end{infobox}

\vspace{1cm}
\noindent{\color{lpAccent}\rule{\linewidth}{0.6pt}}\\[4pt]
{\footnotesize\sffamily\color{lpGray}Lernplatt \textperiodcentered{} by Can Siebert \textperiodcentered{} Kurs 22 (Bo 143) \textperiodcentered{} Tag 7 \textperiodcentered{} Lösungsheft \textperiodcentered{} \textcopyright{} 2026}

\end{document}
